\documentclass[11pt,a4paper]{article}
\usepackage[margin=1in]{geometry}
\usepackage{amsmath,amssymb,amsthm}
\usepackage{algorithm}
\usepackage{algpseudocode}
\usepackage{booktabs}
\usepackage{hyperref}

\theoremstyle{definition}
\newtheorem{definition}{\textbf{Definition}}[section]
\newtheorem{theorem}{\textbf{Theorem}}[section]
\newtheorem{lemma}[theorem]{\textbf{Lemma}}
\newtheorem{remark}{\textbf{Remark}}[section]

\title{\textbf{ALGO-CML-106: CART decision trees}}
\author{\textbf{Team Scaibu} \\ \small Email: \texttt{jaydeep.9664920749@gmail.com} \\ \small \textbf{Architecture and Core Engine Team}}
\date{\textbf{October 2026}}

\begin{document}
\maketitle

\begin{abstract}
Rigorous mathematical formulation, computational complexity derivation, optimization properties, and verifiable execution pipeline for CART decision trees.
\end{abstract}

\section{Definitions and Cognitive Mental Models}

\subsection{Formal Mathematical Definition}
\textbf{Definition (Algorithm \#106 --- CART decision trees):}
Recursive partitioning of the feature space by binary splits, chosen greedily to make each resulting group as pure (classification) or homogeneous (regression) as possible.

Given design observation matrix $\mathbf{X} \in \mathbb{R}^{n \times p}$ and response vector $\mathbf{y} \in \mathbb{R}^n$, the objective operator minimizes statistical empirical risk subject to structural penalty constraints:
\begin{equation}
\hat{\boldsymbol{\beta}} = \arg\min_{\boldsymbol{\beta} \in \mathbb{R}^p} \sum_{i=1}^n \ell(y_i, \mathbf{x}_i^T \boldsymbol{\beta}) + \lambda \Omega(\boldsymbol{\beta})
\end{equation}
where $\ell(\cdot)$ is the loss function and $\Omega(\cdot)$ imposes regularizing structural sparsity.

\subsection{Expanded Technical Specification}
\begin{itemize}
\item \textbf{At a node, for each feature and eac}: h candidate threshold, compute the impurity reduction: Gini impurity or entropy for classification, variance (squared error) for regression.
\item \textbf{Choose the best split. Sorting each}:  feature lets every threshold be evaluated in one scan.
\item \textbf{Recurse on the two children.}: 
\item \textbf{Stop at a maximum depth, a minimum }: number of samples per leaf, or a minimum impurity decrease.
\item \textbf{**Cost-complexity pruning:** grow a}:  large tree, then prune back by a penalty α × (number of leaves), choosing α by cross-validation.
\item \textbf{Leaves predict the majority class o}: r the mean value.
\end{itemize}

\subsection{Computational Complexity and Cost Model}
\begin{equation}
\mathcal{O}(n \cdot p + p^3) \quad \text{operations for standard matrix decomposition and solve}
\end{equation}
\textbf{Operational Resource Cost:} About O(p × n log n) per level.

\section{Algorithmic Pipeline Specification}

\begin{algorithm}[H]
\caption{CART decision trees Execution}
\begin{algorithmic}[1]
\Require Input data $\mathbf{X} \in \mathbb{R}^{n \times p}$, target observations $\mathbf{y} \in \mathbb{R}^n$, regularization $\lambda$
\Ensure Optimized model parameters $\hat{\boldsymbol{\beta}}$, diagnostic telemetry $\mathcal{M}$
\State Center and scale design columns: $\tilde{\mathbf{X}} \gets \text{Standardize}(\mathbf{X})$
\State Compute sample covariance or Gram matrix $\mathbf{G} \gets \tilde{\mathbf{X}}^T \tilde{\mathbf{X}} + \lambda \mathbf{I}_p$
\State Factorize $\mathbf{G} = \mathbf{L} \mathbf{L}^T$ via stable Cholesky decomposition
\State Solve forward substitution $\mathbf{L} \mathbf{z} = \tilde{\mathbf{X}}^T \mathbf{y}$
\State Solve backward substitution $\mathbf{L}^T \hat{\boldsymbol{\beta}} = \mathbf{z}$
\State Compute residual variance $\hat{\sigma}^2 \gets \frac{1}{n - p} \| \mathbf{y} - \tilde{\mathbf{X}} \hat{\boldsymbol{\beta}} \|_2^2$
\State $\mathcal{M} \gets \text{CompileDiagnostics}(\hat{\boldsymbol{\beta}}, \hat{\sigma}^2)$
\State \Return $\{\hat{\boldsymbol{\beta}}, \mathcal{M}\}$
\end{algorithmic}
\end{algorithm}

\section{Theoretical Guarantees: Consistency and Asymptotic Normality}

\begin{theorem}[Consistency and Asymptotic Normality]
Under regular Gauss-Markov conditions with independent, homoscedastic errors $\epsilon_i \sim (0, \sigma^2)$, the estimator $\hat{\boldsymbol{\beta}}$ is unbiased and asymptotically normally distributed:
\begin{equation}
\sqrt{n}(\hat{\boldsymbol{\beta}} - \boldsymbol{\beta}^*) \xrightarrow{d} \mathcal{N}\left(\mathbf{0}, \, \sigma^2 (\mathbf{X}^T \mathbf{X} / n)^{-1}\right)
\end{equation}
achieving the Cramer-Rao lower bound for minimum variance linear unbiased estimators.
\end{theorem}

\begin{proof}
Expanding $\hat{\boldsymbol{\beta}} = (\mathbf{X}^T \mathbf{X})^{-1} \mathbf{X}^T (\mathbf{X} \boldsymbol{\beta}^* + \boldsymbol{\epsilon}) = \boldsymbol{\beta}^* + (\mathbf{X}^T \mathbf{X})^{-1} \mathbf{X}^T \boldsymbol{\epsilon}$.
Taking expectation yields $\mathbb{E}[\hat{\boldsymbol{\beta}}] = \boldsymbol{\beta}^*$. The covariance matrix is $\text{Cov}(\hat{\boldsymbol{\beta}}) = \sigma^2 (\mathbf{X}^T \mathbf{X})^{-1}$. By the Central Limit Theorem applied to the linear combinations of independent error terms, asymptotic normality holds directly.
\end{proof}

\section{Complexity Analysis}
\begin{itemize}
    \item \textbf{Time Complexity:} $\mathcal{O}(n p^2 + p^3)$ for QR or Cholesky decomposition.
    \item \textbf{Space Complexity:} $\mathcal{O}(p^2)$ memory for the Hessian or Gram matrix.
\end{itemize}

\begin{thebibliography}{9}
\bibitem{hastie2009} T.~Hastie, R.~Tibshirani, and J.~Friedman, ``The Elements of Statistical Learning,'' Springer, 2009.
\bibitem{bishop2006} C.~M.~Bishop, ``Pattern Recognition and Machine Learning,'' Springer, 2006.
\end{thebibliography}

\end{document}
